\section{Complete certificate data}\label{app:data}

The following listing contains every line of the certificate used
in Proposition~\ref{prop:certificate}, without omissions.
Cards $0,\ldots,5$ are the six bases \eqref{eq:bases}; the listing
starts at card $6$. The grammar and bit-coded placements are
specified in Section~\ref{sec:certificate}.
A long logical line may wrap visually. The editable file
\texttt{verification/certificate.txt} has one logical line per numbered card.
Parentheses describe inline applications of exactly the same
combination rule, with coordinates in the enclosing line's frame.

\lstinputlisting[breakatwhitespace=true,frame=single]{../verification/certificate.txt}

\section{Complete reduced-envelope certificate}\label{app:data25}

This is the complete decimal certificate of Appendix~\ref{app:route25}.
It begins at card $6$ and ends at card $2015$, after the six bases
specified there. A marker beginning with \texttt{\#} states the next
card number; it is not itself a card. Each other nonempty logical line
is one composition. Every three-digit block, including its leading
zeros, belongs to the decimal grammar. For readable line breaking,
the displayed numeric rows have a bracketed card number and spaces
between blocks. Remove that bracketed prefix and those spaces to
recover each exact numeric line of the literal file; markers and blank
lines are unchanged. This reversible formatting is checked in the
source distribution. Long rows may wrap visually.

\lstinputlisting[breakatwhitespace=true,frame=single]{../verification/supporting/route25/certificate-display.txt}

\section{Complete letter-and-offset certificate}\label{app:data35}

This listing supplies every row $6$ through $615$ of the $35$-move
certificate. Appendix~\ref{app:route35} defines the six base cards,
every letter and translation offset, and the complete reconstruction.
Each line starts with its row number, followed by a nonempty list of
placed earlier cards. Long logical lines may wrap visually.

\lstinputlisting[breakatwhitespace=true,frame=single]{../verification/supporting/prior-35/certificate.txt}

\section{Complete executable verifier}\label{app:checker}

The verifier below reads the literal certificate and evaluates the
set operations described in the paper. From the paper root,
run
\begin{lstlisting}[basicstyle=\ttfamily\small]
python3 verification/verify_certificate.py
\end{lstlisting}
An explicit input path may be supplied with \texttt{--table}.
The wrapper rejects any data bytes with a different checksum and
refuses optimized Python execution, which would disable the
assertions in the checking function. Successful termination means
that every executed check has passed; the semantic implication is
the induction in Proposition~\ref{prop:certificate} and
Lemma~\ref{lem:combination}.

The arrays \texttt{AA}, \texttt{TT}, and \texttt{HH} hold the required
sets, envelopes, and heights of completed numbered cards.
The recursive function \texttt{expr} evaluates an inline expression
in the current line's coordinate frame. A reference uses an earlier
array entry; it never interprets an inline pivot as a translation.
All computations use exact Python integers and finite sets.

\begingroup
\fontencoding{T1}\selectfont
\renewcommand{\ttdefault}{lmtt}
\lstinputlisting[language=Python,breakatwhitespace=false,
 frame=single]{../verification/verify_certificate.py}
\endgroup

\subsection{Independent reconstructions of all three certificates}

The source distribution contains an additional, independently parsed
reconstruction for every dialect, together with tests of the finite
identities used in the article. From the paper root run
\begin{lstlisting}[basicstyle=\ttfamily\small]
python3 verification/supporting/verify_all.py --output-dir ../all-checks
\end{lstlisting}
Choose a new output directory outside the paper directory. The command
runs the primary and independent checks for all three literal tables
and records their full-card agreement, exact identities, and rejection
tests. The support guide gives each program's standalone command.
The $21$-move programs under \texttt{verification/supporting/route21/} additionally
check every inline node and local reply class; they are distinct from
the short verifier printed above. The $25$-move programs under
\texttt{verification/supporting/route25/} use the reduced sets and swap-first anchor
maps. The programs under \texttt{verification/supporting/route35/} check the
letter-and-offset table and its conditional and tactical identities.
All three literal tables and every checking program are included in
the editable source distribution. The directory
\texttt{verification/supporting/prior-35/} contains only a copy of the $35$-move
literal certificate. The scope of the finite reconstruction supplement
is stated in Appendix~\ref{app:formal}.
